% ═══════════════════════════════════════════════════════════════
%  appendices/b-build.tex
%  Приложение Б. Сборка проекта
% ═══════════════════════════════════════════════════════════════

\chapter{Сборка проекта}
\label{app:build}

В этой книге мы используем \textbf{CMake}~--- стандартную систему
сборки для проектов на C. Она кроссплатформенная, гибкая и понятная.

\section{Зависимости}

Перед сборкой убедитесь, что установлены:

\begin{itemize}
    \item \textbf{GCC} или \textbf{Clang}~--- компилятор C.
    \item \textbf{CMake} (версия 3.10 или выше).
    \item \textbf{GTK 4} и \textbf{GLib}~--- библиотеки интерфейса.
    \item \textbf{libudev}~--- идентификация устройств.
    \item \textbf{pkg-config}~--- для поиска библиотек.
\end{itemize}

\subsection{Установка на Debian/Ubuntu}

\begin{minted}[frame=none, bgcolor=gray!5]{bash}
sudo apt update
sudo apt install build-essential cmake pkg-config \
                 libgtk-4-dev libglib2.0-dev libudev-dev
\end{minted}

\subsection{Установка на Fedora}

\begin{minted}[frame=none, bgcolor=gray!5]{bash}
sudo dnf install gcc cmake pkg-config \
                 gtk4-devel glib2-devel libudev-devel
\end{minted}

\subsection{Установка на Arch Linux}

\begin{minted}[frame=none, bgcolor=gray!5]{bash}
sudo pacman -S base-devel cmake pkg-config \
              gtk4 glib2 systemd-libs
\end{minted}

\section{Сборка из исходников}

\subsection{Шаг 1. Клонирование}

\begin{minted}[frame=none, bgcolor=gray!5]{bash}
git clone https://github.com/danayer/dualboard-stm32-oscilloscope.git
cd dualboard-stm32-oscilloscope/pc-app
\end{minted}

\subsection{Шаг 2. Сборка}

\begin{minted}[frame=none, bgcolor=gray!5]{bash}
make
\end{minted}

\subsection{Шаг 3. Запуск}

\begin{minted}[frame=none, bgcolor=gray!5]{bash}
./osc_gen_ui
\end{minted}

\section{Debug и Release}

CMake поддерживает два режима сборки:

\begin{table}[htbp]
\centering
\caption{Сравнение режимов сборки}
\label{tab:build-modes}
\small
\begin{tabularx}{\textwidth}{@{}l L L@{}}
\toprule
\textbf{Параметр} & \textbf{Debug} & \textbf{Release} \\
\midrule
Флаги        & \texttt{-g3 -O0} & \texttt{-O2} \\
Санитайзеры  & ASan + UBSan & нет \\
Отладка      & gdb, valgrind & не нужна \\
Размер       & больше & меньше \\
Скорость     & медленнее & быстрее \\
\bottomrule
\end{tabularx}
\end{table}

\subsection{Debug-сборка}

\begin{minted}[frame=none, bgcolor=gray!5]{bash}
cmake -DCMAKE_BUILD_TYPE=Debug ..
make -j$(nproc)
\end{minted}

\begin{ideabox}
\textbf{Всегда начинайте с Debug.} Санитайзеры (ASan, UBSan)
поймают большинство ошибок памяти \emph{до} того, как они
превратятся в segfault.
\end{ideabox}

\subsection{Release-сборка}

\begin{minted}[frame=none, bgcolor=gray!5]{bash}
cmake -DCMAKE_BUILD_TYPE=Release ..
make -j$(nproc)
\end{minted}

\section{Правила udev для доступа без sudo}

По умолчанию порт \texttt{/dev/ttyACM0} доступен только root.
Создайте файл \texttt{/etc/udev/rules.d/99-stm32.rules}:

\begin{minted}[frame=none, bgcolor=gray!5]{bash}
SUBSYSTEM=="tty", ATTRS{idVendor}=="0483", MODE="0666"
SUBSYSTEM=="tty", ATTRS{idVendor}=="2341", MODE="0666"
\end{minted}

Примените правила:

\begin{minted}[frame=none, bgcolor=gray!5]{bash}
sudo udevadm control --reload-rules
sudo udevadm trigger
\end{minted}

\section{Прошивка STM32}

Прошивка лежит в \texttt{firmware/oscilloscope/} и
\texttt{firmware/generator/}. Соберите через \textbf{STM32CubeIDE}
или \textbf{Makefile}.

\begin{minted}[frame=none, bgcolor=gray!5]{bash}
cd firmware/oscilloscope
make
st-flash write build/oscilloscope.bin 0x08000000
\end{minted}

\section{Сборка книги}

Если хотите собрать саму книгу:

\begin{minted}[frame=none, bgcolor=gray!5]{bash}
cd book
./compile.sh
\end{minted}

Требуется \textbf{XeLaTeX} с пакетами \texttt{minted}, \texttt{tcolorbox},
\texttt{tikz}.

\begin{notebox}
\textbf{Для minted нужен Python и Pygments.}
Установите: \texttt{sudo apt install python3-pygments}.
Также нужен флаг \texttt{--shell-escape} в команде компиляции.
\end{notebox}


